\documentclass[12pt,a4paper]{article}
\usepackage[top=22mm,bottom=24mm,left=27mm,right=27mm]{geometry}
\usepackage{amsmath,amssymb}
\usepackage{xcolor}
\definecolor{link}{HTML}{205973}
\usepackage[colorlinks=true,linkcolor=link,urlcolor=link]{hyperref}
\hypersetup{pdftitle={Strong duality for sampled value-function approximation}}
\setlength{\parindent}{0pt}
\setlength{\parskip}{5pt}
\raggedbottom
\newcommand{\ip}[2]{\langle #1,#2\rangle}
\newcommand{\norm}[1]{\lVert #1\rVert}
\newcommand{\R}{\mathbb R}
\newcommand{\lead}[1]{\par\smallskip\noindent\textbf{#1}\enspace}

\begin{document}

{\LARGE\bfseries Strong duality for sampled\\[2pt] value-function approximation\par}
\medskip
When the slack penalty is at least \(1/(1-\alpha)\), where \(\alpha\) is the
discount factor, the sampled problem has a finite maximum. We show how to
recover an optimal weight, intercept and slacks from any solution of a
finite-dimensional dual problem.

\section*{The model and the two optimization problems}
Let \(\mathcal X\) and \(\mathcal A\) be finite, nonempty sets of states and actions.
Action \(a\) at state \(x\) costs \(g(x,a)\in\R\). The next state is \(y\)
with probability \(p(y\mid x,a)\), where the probabilities are nonnegative
and sum to one over \(y\in\mathcal X\). Future costs are discounted by
\(0<\alpha<1\). A feature map \(\Phi:\mathcal X\to H\) takes values in a real
Hilbert space with inner product \(\ip{\cdot}{\cdot}\) and norm \(\norm{\cdot}\).
We approximate the cost-to-go by \(J_{z,b}(x)=\ip{\Phi(x)}{z}+b\), where
\(z\in H\) is the weight and \(b\in\R\) is the intercept.

Let \(\pi\) be a strictly positive distribution on \(\mathcal X\).
Fix any realization \(x_1,\dots,x_N\) of \(N\ge1\) independent draws from
\(\pi\). A repeated state contributes once per draw. For strictly positive
state relevance weights \(\nu(x)\), write \(w_i=\nu(x_i)/\pi(x_i)\)
and \(g_{ia}=g(x_i,a)\). Define the weighted mean feature \(m\) and the
discounted feature difference \(d_{ia}\) by
\[
  m=\frac1N\sum_i w_i\Phi(x_i),\qquad
  d_{ia}=\Phi(x_i)-\alpha\sum_{y\in\mathcal X}p(y\mid x_i,a)\Phi(y).
\]
Here and below \(i\) ranges over draws and \(a\) over actions.

\lead{Primal problem.} Fix a slack penalty \(\kappa>0\) and a weight penalty
\(\Gamma>0\). Write \(h=1-\alpha\) and \(q=\kappa/N\). Choose \(z\in H\),
\(b\in\R\) and \(s=(s_1,\ldots,s_N)\in\R^N\) to solve
\[
 \begin{aligned}
    \text{maximize}\quad &P(z,b,s)=\ip{m}{z}+b-q\sum_i s_i-\tfrac{\Gamma}{2}\norm{z}^2\\
    \text{subject to}\quad &\ip{d_{ia}}{z}+hb\le g_{ia}+s_i
                               \quad\text{for every }i,a,\\
                          &s_i\ge0 \quad\text{for every }i.
 \end{aligned}
  \tag{P}
\]
The same slack \(s_i\) covers every action at draw \(i\). It bounds how far
\(J_{z,b}(x_i)\) may exceed the immediate cost plus the discounted expected
continuation value.

\lead{Dual problem.} Assign a multiplier \(\lambda_{ia}\) to each draw
and action. Write \(\ell_i(\lambda)=\sum_a\lambda_{ia}\) for the mass in row
\(i\), and \(v(\lambda)=\sum_{i,a}\lambda_{ia}d_{ia}\) for the weighted sum of
feature differences. The required total mass is \(M=1/h\). Solve
\[
 \begin{aligned}
    \text{minimize}\quad &G(\lambda)=\sum_{i,a}g_{ia}\lambda_{ia}
       +\tfrac{1}{2\Gamma}\norm{m-v(\lambda)}^2\\
    \text{subject to}\quad &\lambda_{ia}\ge0,\qquad \ell_i(\lambda)\le q,\qquad
                              \sum_i\ell_i(\lambda)=M .
 \end{aligned}
  \tag{D}
\]
Denote the set of feasible multipliers by \(\Lambda\).

\newpage
\section*{Strong duality and weight recovery}
\textbf{Proposition.} Fix any realization of the sample and assume
\(\kappa\ge 1/(1-\alpha)\). Then (P) and (D) attain the same finite optimal
value. For every minimizer \(\lambda^*\) of \(G\) on \(\Lambda\), the weight
of every primal maximizer is
\[
             z^*=\frac{m-v(\lambda^*)}{\Gamma}.
             \tag{1}\label{eq:weight}
\]
In particular the optimal weight is unique. The worked example illustrates
how the intercept and slacks can vary at the threshold.

\lead{Existence of a dual minimizer.}
The hypothesis gives \(M/N\le\kappa/N=q\). Place mass \(M/N\) on one
action in each row and zero on the other actions. This assignment respects
every row cap and has total mass \(M\), even when \(M=\kappa\).
Thus \(\Lambda\) is nonempty. Its linear constraints define a closed set,
and each multiplier lies between \(0\) and \(q\). Hence \(\Lambda\) is
compact in a finite-dimensional space. The continuous function \(G\)
attains a minimum there.

\lead{Why the dual bounds the primal.}
For a primal point \((z,b,s)\), let
\(\rho_{ia}=g_{ia}+s_i-\ip{d_{ia}}{z}-hb\) be the margin in its
\((i,a)\) constraint. Add these margins to \(P\) with nonnegative weights
\(\lambda_{ia}\). Collecting terms gives
\[
\begin{aligned}
 P+\sum_{i,a}\lambda_{ia}\rho_{ia}
   ={}&\sum_{i,a}g_{ia}\lambda_{ia}
        +\ip{m-v(\lambda)}{z}-\tfrac{\Gamma}{2}\norm{z}^2\\
     &+\Bigl(1-h\sum_i\ell_i(\lambda)\Bigr)b
       +\sum_i\bigl(\ell_i(\lambda)-q\bigr)s_i .
\end{aligned}
\]
For this expression to have a finite supremum over unrestricted \(b\)
and nonnegative \(s\), the coefficient of \(b\) must vanish and each
coefficient of \(s_i\) must be nonpositive. These are exactly the mass
and capacity constraints in (D). For \(\lambda\in\Lambda\), put
\(z_\lambda=(m-v(\lambda))/\Gamma\). Completing the square in \(z\)
gives the exact gap
\begin{equation}
  G(\lambda)-P(z,b,s)
  =\frac{\Gamma}{2}\norm{z-z_\lambda}^2
      +\sum_{i,a}\lambda_{ia}\rho_{ia}
      +\sum_i\bigl(q-\ell_i(\lambda)\bigr)s_i .
  \tag{2}\label{eq:gap}
\end{equation}
For feasible \((z,b,s)\), every term on the right of \eqref{eq:gap} is
nonnegative. Hence every dual value bounds every primal value from above.
Equality forces \(z=z_\lambda\). To finish the proof, we construct an
intercept and slacks that make the gap zero at a dual minimizer.

\newpage
\section*{Constructing a primal optimizer}
Fix any minimizer \(\lambda^*\) of \(G\) and set
\(z=z_{\lambda^*}\). For this fixed weight define the adjusted costs
\(c_{ia}=g_{ia}-\ip{d_{ia}}{z}\). The primal constraints now read
\(hb\le c_{ia}+s_i\). We will obtain \(hb\) from the cheapest action in each row.

\lead{Why the adjusted costs suffice.}
Write \(L(\lambda)=\sum_{i,a}c_{ia}\lambda_{ia}\). Fix any
\(\lambda\in\Lambda\). For \(0<t\le1\) the point
\(\lambda_t=(1-t)\lambda^*+t\lambda\) is feasible. Expanding the square in
the dual objective gives
\[
  0\le\frac{G(\lambda_t)-G(\lambda^*)}{t}
   = L(\lambda)-L(\lambda^*)
     +\frac{t}{2\Gamma}\norm{v(\lambda)-v(\lambda^*)}^2 .
\]
The last term tends to zero as \(t\) decreases to zero. Thus \(\lambda^*\)
minimizes the linear cost \(L\) over \(\Lambda\).

\lead{How to attain the minimum linear cost.}
Let \(r_i=\min_a c_{ia}\). Order the \(N\) rows by \(r_i\), from cheapest
to most expensive. Let \(u\) be the cost of the first row at which the
cumulative capacity reaches or exceeds \(M\). This row exists since \(Nq\ge M\).
The rows strictly cheaper than \(u\) have capacity less than \(M\).
The rows costing at most \(u\) have capacity at least \(M\).

Write \((t)_+=\max(t,0)\). For every feasible allocation,
\[
 L(\lambda)\ \ge\ uM+\sum_i(r_i-u)\ell_i(\lambda)
                  \ \ge\ uM-q\sum_i(u-r_i)_+ .
\]
The first inequality uses \(c_{ia}\ge r_i\). For the second, a row cheaper
than \(u\) holds at most \(q\), and a more expensive row contributes a
nonnegative amount. To attain both bounds, fill each cheaper row to capacity,
split the remaining mass equally among the rows tied at \(u\), and leave
the more expensive rows empty. In each occupied row, put its mass on an action
of cost \(r_i\). At \(M=Nq\) this fills every row. Since \(\lambda^*\)
minimizes \(L\), its cost is exactly this lower bound.

\lead{The intercept and slacks.}
The cutoff \(u\) gives an intercept and slacks through
\[
               b=\frac{u}{h},\qquad s_i=(u-r_i)_+ .
\]
Indeed \(c_{ia}+s_i\ge r_i+(u-r_i)_+\ge u=hb\), so the point is feasible.
The construction gives
\[
      L(\lambda^*)=uM-q\sum_i(u-r_i)_+=b-q\sum_i s_i .
\]
In \eqref{eq:gap} the square vanishes because \(z=z_{\lambda^*}\).
The other two terms sum to \(L(\lambda^*)-b+q\sum_i s_i=0\).
This primal point therefore attains the dual minimum. Applying
\eqref{eq:gap} to any optimal pair proves the weight formula~\eqref{eq:weight}.
\hfill\(\square\)

\newpage
\section*{A worked example at the threshold}
Take two states \(x,y\), one action, and one draw at \(x\). The action at
\(x\) costs \(1\) and moves to \(y\). At \(y\) it costs \(0\) and stays
there. Use scalar features \(\Phi(x)=1,\ \Phi(y)=0\), uniform
\(\nu=\pi\), discount \(\alpha=1/2\), and penalties \(\Gamma=\kappa=2\).
Then \(w_1=1\) and the sole feature difference and the mean feature both
equal \(1\). The dual has one multiplier. Its required mass and its cap
both equal \(2\), so \(\lambda^*=2\). Direct substitution gives
\[
     z^*=\frac{1-2}{2}=-\tfrac12,\qquad
     G(2)=2+\frac{(1-2)^2}{4}=\tfrac94 .
\]
The adjusted cost is \(1-(-1/2)=3/2\), so the construction gives
\(b=3,\ s_1=0\). The primal constraint holds with equality and
\(P(-1/2,3,0)=-1/2+3-1/4=9/4\).
Increasing \(b\) by any \(t\ge0\) and \(s_1\) by \(t/2\) preserves both
feasibility and value. Only the primal weight is forced in this example.

\section*{Differences from the original proposition}
Proposition~1 of Bhat, Farias, Moallemi and Zheng asserts, for positive
penalties, that the primal and their quadratic program attain the same
finite value and that an optimal multiplier recovers an optimal primal
weight. Two changes give the theorem proved here.

\lead{1. The penalty threshold.}
The total row capacity is \(\kappa\), while the required mass is \(1/(1-\alpha)\).
This gives the lower bound in the proposition. To see what can fail, take one
state, one action that returns to that state, and one draw. Set
\(\nu=\pi=1,\ H=\R,\ \Phi=0,\ \alpha=1/2,\ \Gamma=2\).
With cost \(g=0\) and \(\kappa=1\), the feasible points
\((z,b,s_1)=(0,2t,t)\), \(t\ge0\), have primal value \(t\).
The dual requires mass \(2\) in a row of capacity \(1\).

\lead{2. The value of the quadratic objective.}
In the notation of this note, the original quadratic program minimizes
\[
  F(\lambda)=\tfrac12\norm{v(\lambda)}^2-\ip{m}{v(\lambda)}
              +\Gamma\sum_{i,a}g_{ia}\lambda_{ia}+C,
  \qquad C=-N^2\norm{m}^2 ,
\]
over the same set \(\Lambda\). Expanding the square in \(G\) gives
\[
     G(\lambda)=\frac{F(\lambda)-C+\tfrac12\norm{m}^2}{\Gamma}.
\]
Since \(\Gamma>0\), the objectives have exactly the same minimizers, and
the original weight formula applies to every one of them. Their numerical
values can differ. In the one-state example above, change the cost to
\(g=1\) and the penalty to \(\kappa=4\). The only feasible multiplier is
\(2\), with \(F(2)=4\) and \(G(2)=2\). The feasible point \((0,2,0)\)
attains the primal value \(2\). In the two-state example, \(C=-1\), and the
same conversion takes \(F(2)=3\) to \(G(2)=9/4\).

\smallskip
{\small
\textbf{Source.} N.~Bhat, V.~F.~Farias, C.~C.~Moallemi and A.~T.~Zheng,
\href{https://moallemi.com/ciamac/papers/kernel-adp-2012.pdf\#page=8}{\emph{Non-parametric Approximate Dynamic Programming via the Kernel Method}},
August 29, 2022, Proposition~1.\par
\textbf{Formal proof.} \href{https://github.com/nikhilpb/kadp/blob/main/Kadp/PrimalAttainment.lean}{Lean theorem \texttt{strongDuality}}.
}
\end{document}
